\chapter*{Заключение}
\addcontentsline{toc}{chapter}{Заключение}
\label{ch:conclusion}

В ходе выполнения лабораторной работы~\cfgLabNumber{} был развёрнут и настроен
почтовый сервер \textbf{iRedMail} на виртуальной машине \texttt{mail}
под управлением Ubuntu~22.04~LTS.

Выполненные задачи:
\begin{itemize}
  \item подготовлена виртуальная машина \texttt{mail}: назначен статический
        IP \texttt{192.168.\cfgLabStudentN.5/24}, настроен FQDN
        \texttt{mail.yazikov.iks531.local};
  \item добавлены A- и PTR-записи для \texttt{mail} в DNS-сервер BIND9
        (ВМ \texttt{gateway});
  \item выполнена установка iRedMail~1.6.2 с бэкендом \textbf{OpenLDAP}
        и веб-сервером \textbf{Nginx};
  \item настроен почтовый домен \texttt{yazikov.iks531.local};
  \item создан почтовый пользователь \texttt{user1@yazikov.iks531.local}
        через панель \textbf{iRedAdmin};
  \item проверена отправка и получение писем через \textbf{Roundcubemail};
  \item подтверждена доставка письма через лог \texttt{/var/log/mail.log}.
\end{itemize}

Полученные навыки --- развёртывание и базовое администрирование iRedMail,
работа с OpenLDAP, Postfix, Dovecot, управление пользователями через iRedAdmin ---
соответствуют компетенциям администратора корпоративной почтовой инфраструктуры.

\endinput
